\documentclass[11pt]{article}
\usepackage[margin=1in]{geometry}
\usepackage{amsmath,amssymb,amsthm}
\usepackage{listings}
\usepackage{xcolor}
\usepackage{enumitem}
\usepackage[colorlinks=true,linkcolor=blue,urlcolor=blue]{hyperref}

\definecolor{codebg}{rgb}{0.97,0.97,0.97}
\definecolor{codekw}{rgb}{0.10,0.10,0.60}
\definecolor{codecm}{rgb}{0.30,0.55,0.30}
\lstset{
  basicstyle=\ttfamily\footnotesize,
  backgroundcolor=\color{codebg},
  keywordstyle=\color{codekw}\bfseries,
  commentstyle=\color{codecm}\itshape,
  language=Python,
  showstringspaces=false,
  breaklines=true,
  frame=single,
  framesep=4pt,
  columns=fullflexible,
}

\newcommand{\R}{\mathbb{R}}
\newcommand{\Sub}[1]{\mathcal{S}_{#1}}
\newcommand{\inner}[2]{\langle #1,\, #2\rangle}

\title{\S12 Principal Angles between Per-Sample Curvature Subspaces\\
\large How the min / max / mean principal angles are computed in code}
\author{eos\_widget --- \texttt{\_sec12\_payload}}
\date{}

\begin{document}
\maketitle

\section{What is being measured}

For each training sample $i$ the per-sample function Hessian
$Q_i=\nabla^2 f_i(\theta)\in\R^{p\times p}$ (the Hessian of the ground-truth-label
network output $f_i$ with respect to the $p$ parameters) has an eigendecomposition.
We summarise each $Q_i$ by a low-dimensional \emph{curvature subspace}: the span of
its top-$k$ eigenvectors (largest eigenvalues) together with its bottom-$k$
eigenvectors (most negative eigenvalues). The diagnostic then asks: \emph{how aligned
are these curvature subspaces across different samples?} The alignment of two
subspaces is quantified by their \textbf{principal angles}, and Panels~1--3 plot the
$N\times N$ grid of pairwise angles together with its min, max, and mean over sample
pairs, as a function of training step.

\section{The code}

The computation is the function \texttt{pa(k)} together with \texttt{offstats},
reproduced here (ASCII-transliterated; \texttt{TV},\texttt{BV} are the per-sample
top/bottom eigenvectors of shape $(N,K,p)$):

\begin{lstlisting}
offm = ~torch.eye(N, dtype=torch.bool, device=dev)   # off-diagonal (i != j) mask

# principal angles (gauge-free: svdvals invariant to eigenvector sign) for k=1,2,5,10,kfull
def pa(k):
    kk = max(1, min(k, K))
    E  = torch.cat([TV[:, :kk, :], BV[:, :kk, :]], dim=1)        # (N, 2kk, p)
    sv = torch.linalg.svdvals(torch.einsum('iap,jbp->ijab', E, E)).clamp(-1.0, 1.0)
    return (torch.acos(sv) * (180.0 / pi)).mean(dim=2)           # (N, N) mean angle per pair

def offstats(g):                                     # min/max/mean over off-diagonal pairs (i != j)
    s = g[offm]
    return float(s.min()), float(s.max()), float(s.mean())
\end{lstlisting}

We now translate each line into mathematics.

\section{Step 1: the per-sample subspace}

For a chosen rank $k$ let $kk=\min(k,K)$. Stack sample $i$'s top-$kk$ and bottom-$kk$
eigenvectors as the rows of
\begin{equation}
  E_i^{(k)}=
  \begin{bmatrix}
    u_{i,1}^{\top}\\ \vdots\\ u_{i,kk}^{\top}\\
    u_{i,-1}^{\top}\\ \vdots\\ u_{i,-kk}^{\top}
  \end{bmatrix}\in\R^{2kk\times p},
  \qquad
  \texttt{E = cat([TV[:,:kk], BV[:,:kk]])}.
\end{equation}
Each $u$ is a unit vector in parameter space $\R^p$, and within a sample the $2kk$
eigenvectors are mutually orthonormal (they are Ritz vectors of one
fully-reorthogonalised Lanczos basis), so $E_i^{(k)}$ has orthonormal rows. It is
therefore an orthonormal basis for the $2kk$-dimensional subspace
\begin{equation}
  \Sub{i}^{(k)}
  =\operatorname{span}\bigl\{u_{i,1},\dots,u_{i,kk},\,u_{i,-1},\dots,u_{i,-kk}\bigr\}
  \subseteq\R^{p},
  \qquad \dim \Sub{i}^{(k)}=2kk .
\end{equation}
In the code, $k$ ranges over $\{1,2,5,10,k_{\mathrm{full}}\}$ with
$k_{\mathrm{full}}=\min(15,K)$.

\section{Step 2: principal angles between two samples}

For a pair $(i,j)$, the einsum \texttt{'iap,jbp->ijab'} forms the cross-Gram matrix
\begin{equation}
  M_{ij}=E_i^{(k)}\,\bigl(E_j^{(k)}\bigr)^{\top}\in\R^{2kk\times 2kk},
  \qquad
  \bigl(M_{ij}\bigr)_{ab}=\inner{e_{i,a}}{e_{j,b}},
\end{equation}
i.e.\ all pairwise inner products between the eigenvectors of sample $i$ and those of
sample $j$. Let its singular values be
\begin{equation}
  \sigma_1\ge\sigma_2\ge\cdots\ge\sigma_{2kk}, \qquad
  \sigma_\ell=\sigma_\ell(M_{ij})\in[0,1]
  \qquad(\texttt{svdvals}, \text{ clamped to } [-1,1]).
\end{equation}
By the variational (CS / Jordan) characterisation of \textbf{principal angles}
$0\le\theta_1\le\cdots\le\theta_{2kk}\le \tfrac{\pi}{2}$ between two subspaces with
orthonormal bases $U,V$, the singular values of $U^{\top}V$ are exactly the cosines of
those angles:
\begin{equation}
  \cos\theta_\ell=\sigma_\ell,
  \qquad\Longrightarrow\qquad
  \theta^{(k)}_{ij,\ell}=\arccos(\sigma_\ell)\cdot\frac{180}{\pi}\ \text{(degrees)}.
\end{equation}
The clamp to $[-1,1]$ only guards against a singular value numerically exceeding $1$,
which would make $\arccos$ undefined. Recall the recursive meaning: $\cos\theta_1$ is
the largest possible correlation $\inner{u}{v}$ over unit vectors $u\in\Sub{i}$,
$v\in\Sub{j}$; each subsequent angle maximises correlation subject to orthogonality to
the previous principal directions.

\paragraph{Gauge invariance.} Eigenvectors are only defined up to sign (and up to an
orthogonal rotation inside any degenerate eigenspace). Replacing $E_i\mapsto O_i E_i$
for orthogonal $O_i$ leaves the singular values of
$M_{ij}=E_iE_j^{\top}\mapsto O_iE_iE_j^{\top}O_j^{\top}$ unchanged. Hence
$\theta^{(k)}_{ij,\ell}$ depends only on the \emph{subspaces}
$\Sub{i}^{(k)},\Sub{j}^{(k)}$, not on the eigenvector signs or basis choice --- this
is the ``gauge-free'' comment, and is why we use singular values of the cross-Gram
rather than raw dot products.

\section{Step 3: two levels of aggregation}

\paragraph{Inner mean ($\texttt{.mean(dim=2)}$).} The $2kk$ angles of a single pair are
averaged into one scalar, giving the $N\times N$ grid
\begin{equation}
  \bar G^{(k)}_{ij}=\frac{1}{2kk}\sum_{\ell=1}^{2kk}\theta^{(k)}_{ij,\ell}
  \quad\in[0^\circ,90^\circ].
\end{equation}
This $\bar G^{(k)}\in\R^{N\times N}$ is the \textbf{Panel-1 heatmap}; by construction
$\bar G^{(k)}_{ii}=0$ (a subspace makes a zero angle with itself).

\paragraph{Outer min / max / mean ($\texttt{offstats}$).} The mask
$\texttt{offm}=\{(i,j):i\ne j\}$ removes the trivial diagonal, and the three reported
statistics are taken over the $N(N-1)$ off-diagonal ordered pairs:
\begin{align}
  \theta^{(k)}_{\min}  &= \min_{i\ne j}\ \bar G^{(k)}_{ij}, &
  \theta^{(k)}_{\max}  &= \max_{i\ne j}\ \bar G^{(k)}_{ij}, &
  \theta^{(k)}_{\mathrm{mean}} &= \frac{1}{N(N-1)}\sum_{i\ne j}\bar G^{(k)}_{ij}.
\end{align}
These are the fields \texttt{mn}, \texttt{mx}, \texttt{me}; Panels~2--3 plot them
versus the training step (one curve per rank $k$).

\paragraph{Whole-grid mean (\texttt{gm}).} The Panel-1 title uses the mean over
\emph{all} $N^2$ entries, including the $N$ diagonal zeros, so it is just the
off-diagonal mean diluted by the diagonal:
\begin{equation}
  \texttt{gm}^{(k)}=\frac{1}{N^2}\sum_{i,j}\bar G^{(k)}_{ij}
  =\frac{N(N-1)}{N^2}\,\theta^{(k)}_{\mathrm{mean}}
  =\frac{N-1}{N}\,\theta^{(k)}_{\mathrm{mean}}.
\end{equation}

\section{Interpretation}

All quantities are in degrees: $0^\circ$ means two samples share the same curvature
subspace, $90^\circ$ means the subspaces are orthogonal.
\begin{itemize}[nosep]
  \item $\theta^{(k)}_{\min}$ --- the most-aligned pair of samples (smallest angle):
        the strongest shared top/bottom curvature directions in the dataset.
  \item $\theta^{(k)}_{\max}$ --- the most-misaligned pair.
  \item $\theta^{(k)}_{\mathrm{mean}}$ --- the typical pairwise misalignment across
        all sample pairs.
\end{itemize}
Tracking these versus the training step shows whether the per-sample curvature
subspaces become more or less aligned as progressive sharpening / edge-of-stability
proceeds.

\end{document}
